\documentclass[10pt]{article}
\usepackage[letterpaper,margin=0.68in]{geometry}
\usepackage[hidelinks]{hyperref}
\usepackage{enumitem}
\usepackage{titlesec}
\usepackage{tabularx}

\setlist[itemize]{leftmargin=1.2em,itemsep=1.5pt,topsep=2pt}
\titleformat{\section}{\large\bfseries}{}{0em}{}[\titlerule]
\titlespacing*{\section}{0pt}{7pt}{4pt}

\newcommand{\entry}[4]{
  \noindent\begin{tabularx}{\textwidth}{@{}Xr@{}}
    \textbf{#1} & #2 \\
    \textit{#3} & \textit{#4}
  \end{tabularx}\vspace{1pt}
}

\begin{document}

\noindent
{\LARGE \textbf{Piter Z. Garcia Bautista}}\\[4pt]
Rochester, NY \textbar\ +1 (646) 288-3300 \textbar\ \href{mailto:pzg8794@rit.edu}{pzg8794@rit.edu}\\
\href{https://linkedin.com/in/garciapiterz}{linkedin.com/in/garciapiterz} \textbar\
\href{https://github.com/pzg8794}{github.com/pzg8794} \textbar\
\href{https://garciapiterz.com}{garciapiterz.com}
\vspace{6pt}

\section{Research Profile}
Graduate researcher, educator, and NSF Robert Noyce Scholar working at the intersection of artificial intelligence, data science, inclusive computer-science education, and reproducible experimentation. Current work includes fairness-aware sequential decision systems, quantum-network routing, healthcare analytics, and classroom uses of AI that preserve human judgment. More than ten years of engineering and data-systems experience support the translation of research ideas into reliable tools, learning activities, and documented evaluation workflows.

\section{Research Interests}
AI for education; learning-science-informed intelligent systems; fairness-aware machine learning; human-centered evaluation; neurodiversity and accessible computing education; reproducible benchmarking; sequential decision-making; quantum-classical hybrid methods.

\section{Education}

\entry{Rochester Institute of Technology}{Expected Dec 2026}{Master of Science, Data Science --- GPA 3.9}{Rochester, NY}
\begin{itemize}
  \item Concentration: Bioinformatics \& Artificial Intelligence.
  \item Relevant coursework: Foundations of Data Science (DSCI-633), Software Engineering for Data Science (DSCI-644), Applied Data Science I (DSCI-601), Bioinformatics Algorithms (BIOL-630), Data-Driven Knowledge Discovery (ISTE-780).
\end{itemize}

\entry{University of Rochester, Warner School of Education}{Expected Dec 2026}{Master of Science, Teaching Computer Science K--12 --- GPA 4.0}{Rochester, NY}
\begin{itemize}
  \item NSF Robert Noyce Teacher Scholarship Program (full tuition waiver).
  \item Advanced Certificates in Teaching Students with Disabilities, All Grades, and Digitally-Rich Teaching in K--12 Schools (in progress).
\end{itemize}

\entry{Rochester Institute of Technology}{2015}{Master of Science, Computer Science --- GPA 3.2}{Rochester, NY}
\begin{itemize}
  \item Concentrations: Computer Graphics, Artificial Intelligence, Big Data Analytics.
  \item Coursework in compiler design, algorithm analysis, and distributed systems.
\end{itemize}

\entry{Farmingdale State College}{2012}{Bachelor of Science, Computer Engineering Technology --- GPA 3.3}{Farmingdale, NY}
\begin{itemize}
  \item Dual degree: Electrical Engineering \& Computer Engineering Technology.
  \item Honors: CSTEP Award \& Merit, Dual BS Scholarship.
\end{itemize}

\section{Research and Professional Experience}

\entry{Rochester Institute of Technology}{Fall 2025 -- Present}{Graduate Assistant, Quantum and AI Research}{Rochester, NY}
\begin{itemize}
  \item Build a Python multi-testbed framework for adaptive quantum routing and qubit allocation using multi-armed, contextual, adversarial, and neural bandit methods.
  \item Design reusable feedback, validation, logging, dataset, and experiment-tracking workflows across stochastic and adversarial environments.
  \item Evaluate utility, latency, reliability, robustness, and group-level outcomes rather than relying on one aggregate score.
\end{itemize}

\entry{University of Rochester / Pine Brook Elementary School}{Feb 2026 -- May 2026}{Computer Science Student Teacher}{Rochester, NY}
\begin{itemize}
  \item Co-designed and delivered K--5 computer science instruction using computational thinking, inclusive pedagogy, Universal Design for Learning, and multiple participation formats.
  \item Produced lesson artifacts and instructional analyses connecting AI literacy, accessible education, and evidence-based computer-science pedagogy.
\end{itemize}

\entry{Rochester Institute of Technology}{2024 -- 2025}{Equitable Bioinformatics Research --- ISTE-780}{Rochester, NY}
\begin{itemize}
  \item Applied systematic fairness analysis to RNA secondary structure prediction pipelines; demonstrated statistical significance ($p < 0.01$) in bias detection across algorithmic variants using SHAP and disparity-impact analysis.
  \item Developed reproducible evaluation codebase (900+ lines) with automated fairness-metric reporting and statistical testing.
\end{itemize}

\entry{Rochester Institute of Technology}{2025 -- 2026}{Computational Biology Research (BIO614 and BIO550)}{Rochester, NY}
\begin{itemize}
  \item Led formal computational analysis in BIO614: advanced RNA secondary structure prediction with thermodynamic MFE enhancements; produced \textit{BIO614-FinalProjectProposal} (Overleaf manuscript).
  \item Conducted reproducible bulk RNA-seq differential-expression analysis (BIO550): DESeq2 pipeline with QC validation and biological interpretation of nociceptor gene expression.
\end{itemize}

\entry{VEDADATA Inc.}{Sep 2022 -- Aug 2024}{Data Solutions Engineer}{New York, NY}
\begin{itemize}
  \item Designed scalable Python workflows for ingestion, validation, and statistical diagnostics, emphasizing data reliability and structured quality checks.
  \item Strengthened production reliability through anomaly detection, schema validation, and structured logging.
\end{itemize}

\entry{VIOME Inc.}{Jan 2021 -- Sep 2022}{AI Data Solutions Engineer}{New York, NY}
\begin{itemize}
  \item Developed healthcare machine-learning workflows covering preprocessing, feature engineering, supervised experimentation, evaluation, and deployment-oriented data preparation in AWS-based environments.
\end{itemize}

\section{Selected Technical Writing}
\begin{itemize}
  \item \textit{Big Data Medical Diagnosis: A Multi-Method Approach} (RIT M.S.\ Computer Science graduate research, 2015) --- ensemble and deep learning methods for clinical classification on high-dimensional biomedical data.
  \item \textit{Predicting Hospital Readmission Rates: A Multi-Method ML Analysis} (DSCI-633 Project Report, 2025) --- ensemble and regularized regression methods for clinical risk stratification under distribution shift.
  \item \textit{BIO614-FinalProjectProposal} (Overleaf manuscript, 2026) --- RNA secondary structure prediction with thermodynamic deep learning enhancements, reproducible benchmarking, and biological interpretation.
  \item \textit{Equitable Bioinformatics: Enhancing Diagnostic Decision-Making through RNA and Biomarker Data} (ISTE-780 Project Report, 2025) --- fairness-aware ML for genomic diagnostic algorithms with SHAP-based bias detection and statistical significance testing.
  \item \textit{Differential Gene Expression in Murine DRG Neurons Following Sciatic Nerve Injury: An NGS Reanalysis} (BIOL550 Computational Biology Project, 2026) --- reproducible bulk RNA-seq pipeline using DESeq2, QC validation, and pathway-level biological interpretation of nociceptor gene expression.
  \item \textit{Fairness-Aware Bandits for Clinical Decision Systems} (DSCI-601 Project Proposal, 2025) --- bandit learning with equity constraints applied to high-stakes clinical diagnostic decision-making, addressing demographic disparity in sequential clinical recommendation workflows.
  \item \textit{Quantum Network Path Optimization and Fault-Tolerant Routing} (research manuscript, 2025--2026) --- quantum-classical hybrid ML for fault-tolerant routing under noise constraints and decoherence.
\end{itemize}

\section{Awards \& Recognition}
\begin{itemize}
  \item NSF Robert Noyce Teacher Scholarship (2024--2026) --- full tuition waiver for M.S.\ in Teaching Computer Science K--12.
  \item MS International Scholarship (2012--2015) --- MESCYT \& RIT, awarded for academic excellence.
  \item Honorable Mention Technology Award (2012) --- NYS STEP/CSTEP, 2nd place capstone project.
  \item IEEE Alumni Member (2009--Present) --- Member No.\ 90613000.
\end{itemize}

\section{Technical Competencies}
\textbf{Programming:} Python, R, Java, C++, C\#, MATLAB, SQL, JavaScript, Bash\\
\textbf{Research Systems:} Python multi-testbed evaluation, contextual and adversarial bandits, validation, logging, dataset and experiment tracking\\
\textbf{Machine Learning:} PyTorch, TensorFlow, scikit-learn, XGBoost, fairness-aware ML, contextual bandit workflows\\
\textbf{Data \& Analysis:} pandas, NumPy, SciPy, SHAP analysis, statistical testing, reproducible benchmarking\\
\textbf{Infrastructure:} AWS, GCP, Docker, Git/GitHub, Linux, LaTeX

\section{Selected Fit for AI for Education Research}
\begin{itemize}
  \item Combines current AI and sequential-decision research with formal K--12 computer-science teacher preparation and direct elementary, middle-school, and environmental-science teaching experience.
  \item Designs accessible AI-supported learning activities using backward design, Universal Design for Learning, multimodal participation, and explicit routines for questioning and revising model outputs.
  \item Builds reproducible Python-based experimental and evaluation workflows across education, healthcare, bioinformatics, and quantum-network domains.
  \item Brings a disability-informed and cross-cultural perspective to research on how intelligent systems can support learning without replacing educators, peer exchange, or student agency.
\end{itemize}

\end{document}
